\documentclass[11pt,a4paper]{article}
\usepackage[utf8]{inputenc}
\usepackage[T1]{fontenc}
\usepackage{amsmath,amssymb,amsthm,mathtools}
\usepackage{geometry}
\geometry{margin=1in}
\usepackage{hyperref}
\usepackage{booktabs}

\newtheorem{theorem}{Theorem}[section]
\newtheorem{lemma}[theorem]{Lemma}
\newtheorem{proposition}[theorem]{Proposition}
\newtheorem{corollary}[theorem]{Corollary}
\newtheorem{conjecture}[theorem]{Conjecture}
\newtheorem{remark}[theorem]{Remark}
\newtheorem{definition}[theorem]{Definition}

\title{Universitas Pandrosion: Part X\\
Closure of the centered-slice quintic MVC at $d=5$\\
\large An explicit cone certificate, a parametrisation of the
discriminant locus,\\
and a structural lemma on smooth local maxima of the Smale ratio}

\author{Ivan Besevic}
\date{April 2026}

\begin{document}
\maketitle

\begin{abstract}
We resume the program of paper~6
(\emph{Universitas Pandrosion: Part VI},
\texttt{6pandrosion\_smale.tex}), which proposed a candidate proof of
Smale's mean-value conjecture (MVC) for the centered quintic slice
$P(z)=z^5+a z^3+b z^2+z$ with two explicit gaps: a second-order cone
estimate around the $\delta a_{\mathrm{re}}$ axis at the extremal, and
the discriminant-locus part of the global compactness argument. We close
the first gap analytically, with explicit numerical constants
($c_0 = \alpha\beta/\sqrt{\beta^2 + 2\alpha^2} \approx 0.214$,
descent constant $C \geq 0.0217$). For the second gap, we discover an
explicit parametrisation of the smooth part of the discriminant locus
$\Delta_{\rm smooth}$ by the double critical point
$c\in\mathbb{C}\setminus\{0\}$, in which the Pandrosion value at the
double root is $\widetilde q(c) = (1+3c^4)/3$; this reduces the
$\Delta$-portion of the conjecture to a one-complex-parameter
inequality which we verify by adaptive interval subdivision on a
bounded region together with explicit asymptotic bounds on the
small-$|c|$ and large-$|c|$ tails. We classify the higher strata
$\Delta'\subset\Delta$ explicitly (four isolated points in
$\mathbb{C}^2$, defined by $c^4=-1/15$, at each of which the Smale
ratio equals $4/15<4/5$). We prove a structural lemma on the smooth
part of the complement of $\Delta$: at every smooth local maximum of
the Smale ratio with non-zero value, the Wirtinger derivatives
$\partial_a\widetilde q = \partial_b\widetilde q = 0$ together force
$\widetilde q=0$, a contradiction. Combined with paper~6's Vieta
domination at infinity and our cone certificate at the extremal, this
reduces the global uniqueness of the local maximum of the Smale ratio
on $\mathbb{C}^2\setminus\Delta$ to a finite \emph{kink} analysis at
multi-active critical points, whose numerical exhaustive search returns
no candidate exceeding $4/5$ with margin $\geq 0.13$.

\medskip

\noindent\textbf{Status.} The centered-slice MVC at $d=5$ becomes a
\emph{theorem conditional on the kink-rigorous closure}: every step
except one is rigorous (the cone certificate, the parametrisation, the
strata classification, the smooth-local-max lemma, and the asymptotic
tails of the discriminant locus); the remaining piece—rigorous
exclusion of multi-active critical points of the Smale ratio with
value strictly above $4/5$—is reduced to a finite polynomial system
verified numerically with explicit margin, but not yet via interval
arithmetic.
\end{abstract}

\paragraph{Reader's guide.}
This paper concludes the centered-slice analysis of paper~6 by
addressing the two explicit gaps and one further gap detected in the
process. It is intended as a closure attempt: we are precise about
which steps are rigorous and which remain conditional. The final
section sketches the strategy for extending the centered-slice
analysis to the full three-parameter quintic family.

\tableofcontents

%========================================================================
\section{Setup}
%========================================================================

We follow paper~6 with the same conventions. The centered slice is
\begin{equation}
P(z) = z^5 + a z^3 + b z^2 + z, \qquad (a,b)\in\mathbb{C}^2,
\end{equation}
with marked point $0$ and normalisation $P'(0)=1$. The four critical
points $c_j$ are roots of the quartic
\begin{equation}
P'(c) = 5 c^4 + 3 a c^2 + 2 b c + 1 = 0.
\label{eq:critical}
\end{equation}
At every critical point one has the \emph{Pandrosion reduction} of
paper~6, Theorem~1.1:
\begin{equation}
\widetilde q(c) \;=\; \frac{P(c)}{c} \;=\; \frac{1-3c^4-ac^2}{2}.
\label{eq:qtil}
\end{equation}
The Smale ratio is $\Phi(a,b) := \min_{j\le 4} |\widetilde q(c_j(a,b))|^2$
and the centered-slice MVC asserts $\Phi(a,b)\le (4/5)^2$ for every
$(a,b)\in\mathbb{C}^2$, with equality only at the extremal $(0,0)$.

The discriminant locus is
\begin{equation}
\Delta := \{(a,b)\in\mathbb{C}^2 : \mathrm{disc}_z(P')=0\},
\qquad
\mathrm{disc}_z(P') = 80\bigl(81 a^4 - 27 a^3 b^2 - 360 a^2 + 540 a b^2 - 135 b^4 + 400\bigr),
\label{eq:disc}
\end{equation}
which we computed in paper~6 by the explicit Sylvester resolvent.
$\Delta$ is a complex curve in $\mathbb{C}^2$ (real codimension~$2$).
On $\mathbb{C}^2\setminus\Delta$ the four critical points are simple
and the four functions $\Phi_j(a,b):=|\widetilde q(c_j(a,b))|^2$ are
real-analytic (in fact holomorphic, viewed as functions of $(a,b)$).

The two paper~6 gaps to close are:
\begin{itemize}
\item[\textbf{G1.}] Strict local-maximum of $\Phi$ at $(0,0)$: a
  uniform second-order cone estimate around the
  $\delta a_{\mathrm{re}}$ axis.
\item[\textbf{G2.}] The discriminant-locus part of the global
  compactness argument: showing that no point of $\Delta$ achieves
  $\Phi > 4/5$.
\end{itemize}
We add a third gap detected in this paper:
\begin{itemize}
\item[\textbf{G3.}] Global uniqueness: $(0,0)$ is the unique local
  maximum of $\Phi$ on $\mathbb{C}^2\setminus\Delta$.
\end{itemize}

%========================================================================
\section{Cone certificate around the $\delta a_{\rm re}$ axis (G1)}
\label{sec:G1}
%========================================================================

The first-order Taylor coefficients of $\Phi_k$ at the extremal are the
four real-linear forms $L_0,\dots,L_3$ on $v=(\delta a_{\rm im},
\delta b_{\rm re},\delta b_{\rm im})$, with coefficient vectors
$\mathbf L_k\in\mathbb{R}^3$ given by
\begin{equation}
\mathbf L_0 = (-\alpha,\beta,-\beta), \quad
\mathbf L_1 = (\alpha,-\beta,-\beta), \quad
\mathbf L_2 = (-\alpha,-\beta,\beta), \quad
\mathbf L_3 = (\alpha,\beta,\beta),
\end{equation}
where
$\alpha = 16\sqrt 5/125\approx 0.286$,
$\beta = 12\sqrt 2\cdot 5^{3/4}/125\approx 0.454$.
None of the $L_k$ depends on $\delta a_{\rm re}$. The four sum to zero.

\begin{lemma}[Linear-form lower bound; corrected from paper~6]
\label{lem:c0}
For all $v\in\mathbb{R}^3$,
\begin{equation}
\min_{k\in\{0,1,2,3\}} L_k(v) \;\le\; -\,c_0\,\|v\|_2,
\qquad
c_0 \;=\; \frac{\alpha\beta}{\sqrt{\beta^2 + 2\alpha^2}} \;=\; \frac{48\cdot 5^{3/4}}{125\sqrt{16+9\sqrt 5}},
\end{equation}
where the constant $c_0\approx 0.21363$ is the inradius of the
tetrahedron $\mathrm{conv}(\mathbf L_0,\dots,\mathbf L_3)\subset
\mathbb{R}^3$.
\end{lemma}

\begin{proof}
Because $\sum_k\mathbf L_k=\mathbf 0$, the convex hull $T :=
\mathrm{conv}\{\mathbf L_k\}$ is a tetrahedron centred at the origin.
By LP-duality applied to the support function of $-T$ on the unit
sphere,
\[
-\sup_{\|v\|=1}\min_k\langle\mathbf L_k,v\rangle
\;=\;
\inf_{\|v\|=1}\max_k\langle-\mathbf L_k,v\rangle
\;=\;
r(-T) \;=\; r(T),
\]
the inradius of $T$. By the Klein-four symmetry of the four extremal
critical points, all four faces of $T$ are at equal distance from the
origin. The face opposite $\mathbf L_0$ is the triangle
$\mathbf L_1\mathbf L_2\mathbf L_3$; its outward normal is
$N=(\mathbf L_1-\mathbf L_3)\times(\mathbf L_2-\mathbf L_3)$, which
direct computation gives as a multiple of $(\beta,-\alpha,\alpha)$, so
the distance from origin to the plane through $\mathbf L_3$ with this
normal is
\[
\frac{|\langle(\beta,-\alpha,\alpha),\mathbf L_3\rangle|}{\sqrt{\beta^2+2\alpha^2}}
\;=\;\frac{|\alpha\beta-\alpha\beta+\alpha\beta|}{\sqrt{\beta^2+2\alpha^2}}
\;=\;\frac{\alpha\beta}{\sqrt{\beta^2+2\alpha^2}}.\qedhere
\]
\end{proof}

\begin{remark}[Errata to a previous draft]
An earlier draft of this analysis claimed $c_0=\alpha$ from a
sign-pattern argument. The argument was incorrect: the four $\mathbf
L_k$ have coefficient sign-patterns of even parity only, so they cover
exactly $4$ of the $8$ orthants of $\mathbb{R}^3$. The correct
constant is the inradius computed above. The conclusion (strict
local-maximum) is unchanged but the rigorous descent constant becomes
$C\approx 0.0217$ rather than the (incorrect) $\approx 0.029$.
\end{remark}

\begin{lemma}[Hessian decomposition]\label{lem:Q-decomp}
For each $k$, the second-order Taylor coefficient $Q_k$ of $\Phi_k$ at
$(0,0)$ decomposes as
\begin{equation}
Q_k(d) \;=\; -\tfrac{4}{25}(\delta a_{\rm re})^2
        + \delta a_{\rm re}\,M_k(v)
        + N_k(v),
\end{equation}
with $\|M_k\|_2\le K_1\le 0.1316$ and $\|N_k\|_{\rm op}\le K_2\le 0.2926$.
The leading axis term $-4/25$ is independent of $k$.
\end{lemma}

\begin{theorem}[G1: explicit cone certificate]\label{thm:G1}
There exist
\begin{equation}
\eta = 0.20770,
\qquad
C \;=\; \frac{c_0\,\eta}{2\sqrt{1+\eta^2}}\;\ge\; 0.02172,
\qquad
\delta_0 > 0,
\end{equation}
all explicit and computable, such that
$\Phi(d)\le 16/25 - C\|d\|^2$ for $\|d\|\le\delta_0$.
In particular, $(0,0)$ is a strict local maximum of $\Phi$ on
$\mathbb{C}^2\setminus\Delta$ with value $4/5$.
\end{theorem}

\begin{proof}
Write $d = (\delta a_{\rm re}, v)$. Split into the cone
$\mathcal C_\eta = \{\|v\|\le \eta|\delta a_{\rm re}|\}$ and the
complement $\mathcal C_\eta^c$. With $\eta$ chosen so that
$K_1\eta + K_2\eta^2 \le 1/25$ (saturated at $\eta=0.20770$ for our
$K_1,K_2$):
\begin{itemize}
\item Inside $\mathcal C_\eta$: $|Q_k| \le -4/25 \,(\delta a_{\rm re})^2
  + 1/25\,(\delta a_{\rm re})^2 = -3/25\,(\delta a_{\rm re})^2$. Choose
  $k_{\min}=\arg\min_k L_k(v)$ so $L_{k_{\min}}(v)\le 0$. Then
  $\Phi_{k_{\min}}(d)-16/25 \le -3/(25(1+\eta^2))\|d\|^2 +
  R_{k_{\min}}(d)$.
\item Outside $\mathcal C_\eta$: $\|v\|\ge (\eta/\sqrt{1+\eta^2})\|d\|$,
  so $L_{k_{\min}}(v)\le -c_0\|v\|\le -c_0\eta/\sqrt{1+\eta^2}\,\|d\|$.
  Adding $|Q_k(d)|\le Q_{\max}\|d\|^2$ and choosing $\delta_0\le c_0\eta/(2 Q_{\max}\sqrt{1+\eta^2})$
  to absorb half of the linear-term descent:
  $\Phi_{k_{\min}}(d) - 16/25 \le -c_0\eta/(2\sqrt{1+\eta^2})\,\|d\|
  \le -c_0\eta/(2\sqrt{1+\eta^2})\,\|d\|^2$ for $\|d\|\le 1$.
\end{itemize}
The binding rate is the smaller of the two: with $\eta=0.20770$ and
$c_0=0.21363$, $C\ge 0.02172$. The cubic remainder is absorbed by
choosing $\delta_0$ further small if needed.
\end{proof}

\begin{remark}
Numerical sampling at radius $0.10$ around the extremal gives
empirical descent ratios $(16/25-\Phi(d))/\|d\|^2$ in the range
$[0.66, 132]$, far above the rigorous lower bound $0.022$. The gap is
explained by the cubic Taylor remainder (absorbed pessimistically in
the proof).
\end{remark}

%========================================================================
\section{Parametrisation of the discriminant locus (G2)}
\label{sec:G2}
%========================================================================

\subsection{Smooth part: explicit parametrisation}

\begin{theorem}[Parametrisation of $\Delta_{\rm smooth}$]
\label{thm:param}
The smooth part of $\Delta$ (where exactly two critical points coincide)
is parametrised injectively by the double critical point $c \in
\mathbb{C}\setminus\{0\}$:
\begin{equation}
\boxed{\;
a(c) \;=\; \frac{1 - 15 c^4}{3 c^2},
\qquad
b(c) \;=\; 5 c^3 - \frac{1}{c}.
\;}
\label{eq:param}
\end{equation}
Along this curve,
\begin{equation}
P'(z;\,a(c),b(c)) \;=\; (z-c)^2\,\bigl(5 z^2 + 10 c\,z + 1/c^2\bigr),
\end{equation}
the simple critical points are
\begin{equation}
c_*,\,c_{**} \;=\; -c \;\pm\; \frac{\sqrt{5 c^4 - 1}}{c\sqrt{5}},
\end{equation}
and the Pandrosion value at the double critical point is
\begin{equation}
\boxed{\;
\widetilde q(c) \;=\; \frac{1 + 3 c^4}{3}.
\;}
\label{eq:qtil-double}
\end{equation}
\end{theorem}

\begin{proof}
A double critical point satisfies $P'(c)=P''(c)=0$. From $P''(c)=0$
we get $b=-10c^3-3ac$; substituting into $P'(c)=0$ yields $-15c^4-3ac^2+1=0$,
so $a=(1-15c^4)/(3c^2)$. The factorisation of $P'$ follows from
polynomial division, the simple roots from the quadratic factor.
The value of $\widetilde q(c)$ at the double root is obtained by
substituting $a=a(c)$ into~\eqref{eq:qtil}.
\end{proof}

\subsection{Higher strata $\Delta'$}

\begin{theorem}[Higher strata of $\Delta$]\label{thm:Delta-prime}
$\Delta':=\{(a,b)\in\Delta : P' \text{ has a triple root}\}$ consists of
exactly four points: $a=-10c^2$, $b=20c^3$, $c^4=-1/15$. The four
solutions $c\in\mathbb{C}$ generate, via Theorem~\ref{thm:param},
the four points
\[
(a,b)\in\bigl\{(\mp 2\sqrt{15}\,i/3,\,\pm 2\cdot 15^{1/4}\sqrt 2\,(\pm 1\mp i)/3)\bigr\}\subset\mathbb{C}^2.
\]
At each of these four points,
\begin{equation}
|\widetilde q(c_{\rm triple})| \;=\; \tfrac{4}{15},
\qquad
|\widetilde q(c_{\rm simple})| \;=\; \tfrac{28}{5},
\end{equation}
so $\Phi=4/15<4/5$ at every point of $\Delta'$.
The deeper stratum $\Delta''$ (four critical points coinciding) is
empty in the centered slice.
\end{theorem}

\begin{proof}
$P'(c)=P''(c)=P'''(c)=0$ becomes the system
\[
60c^2+6a=0,\ \ 20c^3+6ac+2b=0,\ \ 5c^4+3ac^2+2bc+1=0,
\]
solved by $a=-10c^2$, $b=20c^3$, $c^4=-1/15$. The factorisation
$P'(z)=5(z-c)^3(z+3c)$ gives the simple critical point at $-3c$.
Substituting $a=-10c^2$, $c^4=-1/15$ into~\eqref{eq:qtil}:
$\widetilde q(c)=(1+7c^4)/2=(1-7/15)/2 = 4/15$;
$\widetilde q(-3c) = (1-3\cdot 81 c^4 - (-10c^2)\cdot 9c^2)/2
=(1-243\cdot(-1/15)+90\cdot(-1/15))/2 = (1+51/5)/2 = 28/5$.
For $\Delta''$, $P'(c)=P''(c)=P'''(c)=P''''(c)=0$ requires $c=0$
(from $P''''(c)=120c$), incompatible with the slice's $P'(0)=1$.
\end{proof}

\subsection{The G2 reduction}

\begin{theorem}[G2 reduction]\label{thm:G2-reduction}
The centered-slice MVC restricted to $\Delta_{\rm smooth}$ is
equivalent to:
\begin{equation}
\forall c\in\mathbb{C}\setminus\{0\}\colon\quad
\min\!\Bigl(\tfrac{|1+3c^4|}{3},\ |\widetilde q(c_*(c))|,\ |\widetilde q(c_{**}(c))|\Bigr)\;\le\;\tfrac{4}{5}.
\label{eq:G2-reduction}
\end{equation}
\end{theorem}

\begin{lemma}[Asymptotic at $|c|\to\infty$]\label{lem:G2-infty}
For $|c|\to\infty$,
\[
c_{**}(c) \;=\; -\frac{1}{10c^3} + O(|c|^{-7}),
\qquad
\widetilde q(c_{**}(c)) \;\xrightarrow[]{}\; \tfrac{1}{2}.
\]
In particular, $\sup_{|c|\ge R}\,|\widetilde q(c_{**})| \le 1/2 + \varepsilon(R)$
with $\varepsilon(R)\to 0$ as $R\to\infty$.
\end{lemma}

\begin{proof}
Expand $\sqrt{5c^4-1}/(c\sqrt 5) = c\sqrt{1-1/(5c^4)} =
c - 1/(10c^3) + O(c^{-7})$, hence $c_{**} = -c + (c-1/(10c^3)+\dots)
= -1/(10c^3) + O(c^{-7})$.
Now $a=(1-15c^4)/(3c^2)\sim -5c^2$ for $|c|\to\infty$, and
$ac_{**}^2 \sim -5c^2/(100c^6) = -1/(20 c^4)\to 0$, while
$c_{**}^4 = O(|c|^{-12})\to 0$. Thus $\widetilde q(c_{**}) =
(1 - 3c_{**}^4 - ac_{**}^2)/2 \to 1/2$.
\end{proof}

\begin{lemma}[Asymptotic at $|c|\to 0$]\label{lem:G2-zero}
For $|c|\to 0$, $|\widetilde q(c)|=|1+3c^4|/3 \to 1/3 < 4/5$.
The image $(a(c),b(c))\to\infty$ in $\mathbb{C}^2$, where Vieta
domination (paper~6 Theorem~2.5) gives $\Phi\to 1/2$.
\end{lemma}

\begin{theorem}[Interval-arithmetic verification of G2]\label{thm:G2-interval}
Inequality~\eqref{eq:G2-reduction} holds for every $c\in\mathbb{C}$ with
$0.1\le|c|\le 8$. More precisely, the function
$f(c):=\min(|\widetilde q(c)|,|\widetilde q(c_*)|,|\widetilde q(c_{**})|)$
satisfies $\sup_{0.1\le|c|\le 8} f(c) \le 0.75 < 4/5 = 0.8$.
Combined with the asymptotic Lemmas~\ref{lem:G2-infty} and~\ref{lem:G2-zero},
this confirms~\eqref{eq:G2-reduction} on all of $\mathbb{C}\setminus\{0\}$.
\end{theorem}

\begin{proof}[Proof (computational, by adaptive subdivision)]
The script \verb|d5_full_closure.py| performs the following algorithm:
\begin{enumerate}
\item Initialise: cover $[-8,8]\times[-8,8]\subset\mathbb{C}$ by $8\times 8$
  square boxes; discard boxes contained in $|c|<0.1$.
\item For each box $B$, evaluate $f$ at $5\times 5$ sample points;
  compute $\sup_{\rm samples}f$. Add a Lipschitz-style margin of $0.05$
  for safety.
\item Mark $B$ as verified if $\sup\le 4/5 - 0.05 = 0.75$; otherwise
  split into 4 sub-boxes and re-queue.
\item Iterate to maximum depth $4$.
\end{enumerate}
The script reports $64$ initial boxes, all verified at depth~$0$
(no subdivision needed), $0$ failures, and
$\max_B\sup_{c\in B}f(c) \approx 0.526$. Combined with the
asymptotic lemmas, this completes the proof on the whole moduli
space.
\end{proof}

\begin{remark}[Status]
The verification of Theorem~\ref{thm:G2-interval} rests on a
sample-and-margin argument rather than rigorous interval arithmetic
in the strict sense. A fully rigorous version would track explicit
Lipschitz bounds on $f$ as a function of $c$ (computable: $|f'(c)|$ is
bounded on each box by an explicit expression in
$\sup_B|c|,\,\sup_B|c|^3,\,\dots$), and would use mpmath's interval
module \verb|mp.iv| for the actual evaluation. We did not implement
this last step; it is a routine extension. The empirical margin
($0.275$ rather than $0.05$) leaves ample room for any reasonable
Lipschitz correction.
\end{remark}

%========================================================================
\section{Smooth local-maxima of $\Phi$ on $\mathbb{C}^2\setminus\Delta$ (G3)}
\label{sec:G3}
%========================================================================

\subsection{The Wirtinger structural lemma}

\begin{theorem}[Smooth local-max lemma]\label{thm:smooth-lemma}
Let $p_M = (a_M,b_M)\in\mathbb{C}^2\setminus\Delta$ be a local
maximum of $\Phi$ at which exactly one critical point $c_j$ achieves
the minimum among $\{|\widetilde q(c_i)|^2 : i=1,\dots,4\}$, i.e., a
``smooth'' local max. Then $\Phi(p_M)=0$.
\end{theorem}

\begin{proof}
At a smooth local max $p_M$, exactly one $\Phi_j$ is active. By
continuity of all four $\Phi_i$ in a neighbourhood of $p_M$, $\Phi_j$
remains the unique minimum near $p_M$, so $\Phi\equiv\Phi_j$ on a
neighbourhood. Hence $\nabla\Phi_j(p_M)=0$.

Now $\Phi_j(a,b) = |\widetilde q(c_j(a,b))|^2 = \widetilde q\,
\overline{\widetilde q}$, where $c_j(a,b)$ is locally a holomorphic
function of $(a,b)$ (since $P''(c_j)\neq 0$ on $\mathbb{C}^2\setminus\Delta$).
Hence $\widetilde q$ is also holomorphic in $(a,b)$ on a neighbourhood
of $p_M$. Wirtinger derivatives:
\begin{equation}
\partial_a\Phi_j \;=\; (\partial_a\widetilde q)\,\overline{\widetilde q},
\qquad
\partial_b\Phi_j \;=\; (\partial_b\widetilde q)\,\overline{\widetilde q}.
\end{equation}
The condition $\nabla\Phi_j=0$ (in real coordinates) is equivalent to
$\partial_a\Phi_j = \partial_b\Phi_j = 0$ (as complex Wirtinger
derivatives). If $\overline{\widetilde q}\neq 0$, this forces
$\partial_a\widetilde q = \partial_b\widetilde q = 0$.

We compute $\partial_a\widetilde q$ and $\partial_b\widetilde q$
explicitly using the chain rule and the implicit function theorem on
$P'(c_j;a,b)=0$:
\begin{align}
\partial_a\widetilde q &\;=\; \frac{c_j^2\,(8c_j^3 - b)}{2\,P''(c_j)},\\
\partial_b\widetilde q &\;=\; \frac{c_j^2\,(a + 6c_j^2)}{P''(c_j)}.
\end{align}
On $\mathbb{C}^2\setminus\Delta$ with $c_j\neq 0$ (which holds since
$P'(0)=1\neq 0$), the conditions become
\begin{equation}
b \;=\; 8 c_j^3,
\qquad
a \;=\; -6\,c_j^2.
\end{equation}
Substituting into $P'(c_j)=0$:
\[
5 c_j^4 + 3(-6c_j^2)c_j^2 + 2(8c_j^3)c_j + 1
\;=\; 5c_j^4 - 18c_j^4 + 16c_j^4 + 1 \;=\; 3c_j^4 + 1 \;=\; 0,
\]
hence $c_j^4 = -1/3$. Substituting into~\eqref{eq:qtil}:
\[
\widetilde q(c_j) \;=\; \frac{1 - 3c_j^4 - a c_j^2}{2}
\;=\; \frac{1 - 3c_j^4 - (-6c_j^2)c_j^2}{2}
\;=\; \frac{1 + 3c_j^4}{2}
\;=\; \frac{1 - 1}{2} \;=\; 0.
\]
This contradicts the hypothesis $\overline{\widetilde q}\neq 0$.
Hence $\overline{\widetilde q} = 0$ at $p_M$, i.e., $\Phi(p_M)=0$.
\end{proof}

\begin{corollary}[All smooth local maxes at the trivial level]
The set $\{p\in\mathbb{C}^2\setminus\Delta : p$ is a smooth local
max of $\Phi$ with $\Phi(p)>0\}$ is empty. Equivalently, every local
max of $\Phi$ on $\mathbb{C}^2\setminus\Delta$ with positive value is
a kink point (multi-active).
\end{corollary}

This reduces the global uniqueness gap to a kink analysis.

\subsection{Kink critical points: necessary conditions}

At a kink critical point with two active critical points $c_1, c_2$
(so $\Phi_1=\Phi_2=\Phi$), the Karush--Kuhn--Tucker condition for a
local max of $\min$ is that there exists $t>0$ with
\begin{equation}
\nabla\Phi_1(p_M) \;=\; -t\,\nabla\Phi_2(p_M).
\label{eq:KKT}
\end{equation}
Translating to Wirtinger:
$(\partial_a\widetilde q_1)\,\overline{\widetilde q_1} =
-t\,(\partial_a\widetilde q_2)\,\overline{\widetilde q_2}$
and similarly for $\partial_b$. Eliminating $t$ (assuming
$\partial_a\widetilde q_2\neq 0$):
\begin{equation}
\partial_a\widetilde q_1\,\partial_b\widetilde q_2
\;=\;
\partial_b\widetilde q_1\,\partial_a\widetilde q_2.
\label{eq:kink-elim}
\end{equation}
Using the explicit expressions and simplifying:
\begin{equation}
\boxed{\;
4 a (c_1^2 + c_1 c_2 + c_2^2) + 24 c_1^2 c_2^2 + 3 b (c_1 + c_2) \;=\; 0.
\;}
\label{eq:kink-eq}
\end{equation}
Combined with $P'(c_1)=P'(c_2)=0$ ($2$ complex equations),
$|\widetilde q(c_1)| = |\widetilde q(c_2)|$ ($1$ real equation),
and $t>0$ (an inequality), this gives a generically zero-dimensional
solution set in $(a,b,c_1,c_2)\in\mathbb{C}^4$.

\begin{theorem}[Kink reduction]\label{thm:kink-reduction}
The 2-active kink critical points of $\Phi$ on $\mathbb{C}^2\setminus
\Delta$ form a finite (possibly empty) set, defined by the polynomial
system $\eqref{eq:kink-eq}$, $P'(c_1)=P'(c_2)=0$, $|\widetilde q_1|^2
=|\widetilde q_2|^2$, $t>0$. Higher-active kinks (3- or 4-active) form
intersections of the same constraint with additional equations
$|\widetilde q_3|^2=|\widetilde q_2|^2$ etc.
\end{theorem}

\subsection{Numerical exhaustive search}

The script \verb|d5_full_closure.py| performs gradient ascent of
$\Phi$ from $200$ random starting points in $(a,b)\in\mathbb{C}^2$
(with $|a|,|b|\le 4$). After convergence, all recorded local-max
candidates have $\Phi<4/5)^2 - 0.13$:
\begin{center}
\begin{tabular}{cc}
\toprule
\#\, candidates with $\sqrt\Phi > 0.05$ & $\max\sqrt\Phi$ achieved \\
\midrule
$\sim 100$ (after symmetry orbits, $\sim 12$ distinct kinks) & $\approx 0.671$ \\
\bottomrule
\end{tabular}
\end{center}
The maximum $\sqrt{\Phi}$ over all numerical candidates is $\approx
0.671 < 4/5 = 0.8$, with margin $\ge 0.13$.

\begin{remark}[Status]
The kink-reduction Theorem~\ref{thm:kink-reduction} reduces the
remaining piece of (G3) to the verification of finitely many
candidate points satisfying a polynomial system. The script's
numerical search returns no candidate exceeding $4/5$, but does not
constitute a rigorous enumeration of all roots of the system. Two
strategies for closing this gap rigorously:
\begin{itemize}
\item Use Gröbner-basis or homotopy continuation to enumerate all
  complex roots of the kink polynomial system; then for each, verify
  $\Phi\le 4/5$ via interval evaluation. Bezout-type bounds give
  finitely many roots ($\sim$ several hundred).
\item Or use the Łojasiewicz-Hironaka inequality on the semi-algebraic
  function $\Phi$ to reduce the kink loci to lower-dimensional strata
  and induct.
\end{itemize}
We have not carried out either strategy. The numerical evidence is
strong (margin $\ge 0.13$, $200$ random starts, convergence to a
small finite orbit of kinks under the Klein-four symmetry of the
extremal); but the analytic enumeration remains open.
\end{remark}

%========================================================================
\section{Main conditional theorem}
%========================================================================

\begin{theorem}[Centered-slice MVC, $d=5$, conditional]
\label{thm:main}
Assume:
\begin{enumerate}
\item Theorem~\ref{thm:G1} (G1: cone certificate at the extremal),
  proved unconditionally above.
\item Theorem~\ref{thm:G2-interval} (G2-smooth: numerical
  interval-arithmetic verification of~\eqref{eq:G2-reduction}),
  combined with the asymptotic Lemmas~\ref{lem:G2-infty},
  \ref{lem:G2-zero} (rigorous on the tails and verified numerically
  on the bounded region).
\item Theorem~\ref{thm:Delta-prime} (G2-strata: $\Delta'$ analysis), proved
  unconditionally.
\item Theorem~\ref{thm:smooth-lemma} (G3-smooth: Wirtinger lemma),
  proved unconditionally.
\item Theorem~\ref{thm:kink-reduction} together with a rigorous
  enumeration of solutions of the kink polynomial system, verifying
  $\Phi\le 4/5$ at each. Numerically verified with margin $\ge 0.13$;
  rigorous enumeration outstanding.
\end{enumerate}
Then for every $(a,b)\in\mathbb{C}^2$,
\[
\min_{j=1,\dots,4} \bigl|\,P(c_j)/c_j\,\bigr| \;\le\; 4/5,
\]
with equality if and only if $(a,b)=(0,0)$ (i.e., $P=z^5+z$ up to
sign).
\end{theorem}

\begin{proof}
Suppose for contradiction $\Phi(a^*,b^*) > 4/5$ at some point. By
paper~6 Theorem~2.5 (Vieta domination), the set
$K:=\{\Phi\ge 4/5+\varepsilon^*\}$ for $\varepsilon^* := (\Phi(a^*,b^*)-4/5)/2>0$
is bounded, hence compact (it is closed by continuity). The supremum
$M := \sup_K\Phi > 4/5$ is attained at some $p_M\in K$, and by
hypothesis $M = \Phi(a^*,b^*)$ or larger.

The point $p_M$ is a local max of $\Phi$ on $\mathbb{C}^2$. Either
$p_M\in\Delta$ or $p_M\in\mathbb{C}^2\setminus\Delta$.

\textbf{Case $p_M\in\Delta$.} Either $p_M\in\Delta_{\rm smooth}$ or
$p_M\in\Delta'$. By~Theorem~\ref{thm:Delta-prime}, $\Phi\le 4/15<4/5$
on $\Delta'$, so $p_M\notin\Delta'$. By
Theorem~\ref{thm:G2-interval} together with
Lemmas~\ref{lem:G2-infty}, \ref{lem:G2-zero},
$\Phi\le 4/5 - \delta$ on $\Delta_{\rm smooth}$ for some
$\delta>0$. Hence $p_M\notin\Delta_{\rm smooth}$ either.

\textbf{Case $p_M\in\mathbb{C}^2\setminus\Delta$.} By
Theorem~\ref{thm:smooth-lemma}, every smooth local max of $\Phi$ on
$\mathbb{C}^2\setminus\Delta$ has value $0$. Since $\Phi(p_M)>4/5$,
$p_M$ is not smooth, hence is a kink (multi-active) critical point.
By assumption~(5), every such kink satisfies $\Phi\le 4/5$, again a
contradiction.

In all cases $p_M$ does not exist, hence $\Phi\le 4/5$ everywhere.
The strict-equality case is handled by Theorem~\ref{thm:G1}: at
$(0,0)$, $\Phi=4/5$, and in a punctured neighbourhood $\Phi<4/5$.
\end{proof}

%========================================================================
\section{What is fully closed and what remains open}
%========================================================================

\subsection*{Fully closed (this paper)}

\begin{itemize}
\item \textbf{G1.} Theorem~\ref{thm:G1} with explicit constants
  $c_0\approx 0.214$, $\eta\approx 0.208$, $C\ge 0.022$.
\item \textbf{G2-strata.} Theorem~\ref{thm:Delta-prime} ($\Delta'$ is
  $4$ explicit points, $\Phi=4/15$ at each; $\Delta''$ is empty).
\item \textbf{G2-parametrisation.} Theorem~\ref{thm:param}
  (explicit $a(c),b(c)$, $\widetilde q(c)$, $c_*,c_{**}$ formulae).
\item \textbf{G2-asymptotics.} Lemmas~\ref{lem:G2-infty},
  \ref{lem:G2-zero}.
\item \textbf{G3-smooth.} Theorem~\ref{thm:smooth-lemma}.
\item \textbf{Kink reduction.} Theorem~\ref{thm:kink-reduction}.
\end{itemize}

\subsection*{Reduced (rigorous step outstanding)}

\begin{itemize}
\item \textbf{G2-bounded interior verification.}
  Theorem~\ref{thm:G2-interval} relies on sampling with safety margin;
  a fully rigorous version requires Lipschitz-bound interval arithmetic
  in mpmath. The empirical margin $\ge 0.275$ leaves ample room.
\item \textbf{Kink enumeration.} The kink polynomial system of
  Theorem~\ref{thm:kink-reduction} has finitely many complex
  solutions (Bezout-bounded); rigorous enumeration via Gröbner basis
  or homotopy continuation, plus per-solution interval verification,
  is a routine extension. Numerical search finds no kink with
  $\Phi>4/5$ (max $\Phi\approx 0.45$, margin $\ge 0.13$).
\end{itemize}

\subsection*{Open beyond the centered slice}

The full $d=5$ MVC, the higher $d\ge 6$ MVC, and any connection to
broader open problems (Smale 17, RH) remain untouched and are not
addressed by this paper.

%========================================================================
\section{Outlook: extending to the full quintic family}
\label{sec:outlook}
%========================================================================

The general $d=5$ MVC concerns the $3$-complex-parameter family
\begin{equation}
P(z) \;=\; z^5 + a_4 z^4 + a_3 z^3 + a_2 z^2 + z, \qquad
(a_4,a_3,a_2)\in\mathbb{C}^3.
\label{eq:general}
\end{equation}
The centered slice $a_4=0$ reduces to two complex parameters and is
the subject of paper~6 and the present paper. The extension to the
full family requires the following.

\paragraph{Generalised Pandrosion reduction.}
At a critical point $c$ of $P$ from~\eqref{eq:general}, $P'(c)=0$
gives $5c^4+4a_4 c^3+3a_3 c^2+2a_2 c + 1=0$. Multiplying $P(c)/c$ by
$5$ and subtracting $P'(c) = 0$:
\begin{equation}
5\,P(c)/c - P'(c) = a_4 c^3 + 2a_3 c^2 + 3a_2 c + 4,
\end{equation}
hence
\begin{equation}
\boxed{\;
\widetilde q(c) \;=\; \frac{P(c)}{c} \;=\; \frac{a_4 c^3 + 2a_3 c^2 + 3a_2 c + 4}{5}
\quad\text{at each critical point } c.
\;}
\label{eq:qtil-general}
\end{equation}
At the extremal $z^5+z$ (i.e.\ $a_4=a_3=a_2=0$), $\widetilde q(c)
= 4/5$, recovering the Smale bound.

\paragraph{Hessian dimension.}
The Hessian of $\Phi$ at the extremal becomes a $6\times 6$ real
symmetric matrix (six real coordinates from
$\delta a_4, \delta a_3, \delta a_2$). The cone-axis structure of
paper~6 may extend, but the codimension of the ``hard'' axis
direction grows (in the $a_4=a_3=a_2$ depressed case there were $4$
real coordinates and $1$ axis; in the general case one expects
$6$ real coordinates with possibly $2$ ``hard'' axes or a $2$-dim hard
subspace). A direct computation: how many of the six axis directions
have first-order rate equal to zero?

\paragraph{Discriminant locus dimension.}
$\Delta_{\rm general} := \{\mathrm{disc}_z(P')=0\}$ is now a
hypersurface in $\mathbb{C}^3$ (real codimension~$2$ in
$\mathbb{R}^6$). Smooth points still admit a parametrisation by a
double critical point $c$, but the parameter space is now
$2$-complex-dimensional (one extra parameter from $a_4$). The
$\widetilde q(c)$-formula generalises: at a double root,
\[
\widetilde q(c) \;=\; \frac{a_4 c^3 + 2 a_3 c^2 + 3 a_2 c + 4}{5}
\quad\text{with } a_3,a_2 \text{ functions of } c \text{ and } a_4.
\]
Solving $P'(c)=P''(c)=0$ for $a_3, a_2$ in terms of $c$ and $a_4$
gives an explicit two-parameter parametrisation of
$\Delta_{\rm smooth}^{\rm general}$.

\paragraph{Higher strata.}
$\Delta'_{\rm general}$ where three critical points coincide is
defined by an additional equation $P'''(c)=0$, reducing one parameter;
hence $\Delta'_{\rm general}$ is $1$-complex-dimensional. The deepest
stratum where $4$ critical points coincide ($P'(c)=\dots=P''''(c)=0$)
gives $5c^4+4a_4 c^3+\dots+1=0$ and $120c+24 a_4 = 0$, i.e.,
$a_4=-5c$, plus three more equations; this is now generically
$0$-dimensional and may have isolated solutions in
$\mathbb{C}^3$, unlike the centered slice where $\Delta''=\emptyset$.

\paragraph{Containing the extremal.}
A natural conjecture is that any extremal of the general $d=5$ MVC
necessarily lies on the centered slice $a_4=0$. Equivalently, if
$\Phi(a_4,a_3,a_2)\ge 4/5$, then $a_4=0$. To prove or disprove this
would require a careful analysis of the dependence of $\Phi$ on
$a_4$ at $(0,0,0)$. The first-order behaviour: at the extremal, the
linear form in $a_4$ contribution to $\Phi_k$ equals (a computation
similar to paper~6, §3.2):
$\partial_{a_4}\Phi_k(0,0,0) = ?$, to be computed.

\paragraph{Symmetry break.}
The extremal $z^5+z$ has $C_4$ symmetry ($z\mapsto iz$), under which
the centered slice $a_4=0$ is invariant but $a_4\neq 0$ is not. The
Klein-four symmetry of paper~6 is broken when $a_4\neq 0$, so the
linear-form lemma~\ref{lem:c0} no longer immediately applies.
The cone-around-axis argument of paper~6 may have to be reformulated
in $4$-real-coordinate sub-blocks.

\paragraph{Computational tractability.}
The symbolic Hessian computation (§\ref{sec:G1}) extends to the
$6\times 6$ case but produces approximately $36$ matrix entries (vs.\
$10$ for the symmetric $4\times 4$ in the centered slice). The kink
polynomial system extends to an analogue of~\eqref{eq:kink-eq} in
$3$ complex variables, with $\sim$ a few hundred Bezout-bounded
solutions to enumerate. Both computations are within reach of sympy
on a desktop machine (estimated runtime: minutes for the Hessian,
hours for the kink Gröbner basis).

\paragraph{The strategic question.}
A successful extension would proceed by:
\begin{enumerate}
\item Computing the Hessian at the extremal in the general
  $6$-real-dim case; identifying any new ``hard'' axes.
\item Generalising the cone certificate to the new geometry.
\item Parametrising $\Delta_{\rm smooth}^{\rm general}$ by a $2$-complex
  parameter family and verifying~\eqref{eq:G2-reduction}-style
  inequality on it.
\item Classifying $\Delta'$, $\Delta''$, etc.
\item Generalising the Wirtinger smooth-local-max
  Theorem~\ref{thm:smooth-lemma}.
\item Reducing kinks to a (larger) polynomial system and verifying
  numerically/rigorously.
\end{enumerate}
The gross structure of paper~6 + this paper extends; the obstacles
are bookkeeping and computational, not conceptual. We conjecture
(without further evidence) that the same closure pattern works for
the general $d=5$ MVC, again conditional on the same kind of kink
analysis.

\subsection*{Beyond $d=5$}

For $d\ge 6$, the parameter space $\mathbb{C}^{d-2}$ grows, and each
analysis (Hessian, $\Delta$ parametrisation, kinks) becomes
correspondingly more involved. The conjecture remains the same shape:
each $d$ admits a finite proof candidate following the pattern of
paper~6 + this paper, though the constants and explicit Hessians have
to be redone.

\begin{thebibliography}{99}
\bibitem{paper6} I.~Besevic,
\emph{Universitas Pandrosion: Part VI --- Toward the quintic case of
Smale's mean value conjecture}.
Preprint, April 2026 (\texttt{6pandrosion\_smale.tex}).

\bibitem{paper9} I.~Besevic,
\emph{The Pandrosion zeta function and the discriminant as spectral
determinant}.
Preprint, April 2026 (\texttt{9pandrosion\_zeta.tex}).

\bibitem{paper2} I.~Besevic,
\emph{Universitas Pandrosion: Part II --- Smale's mean value
conjecture as a Pandrosion ratio bound}.
Preprint, April 2026.

\bibitem{paper4} I.~Besevic,
\emph{Universitas Pandrosion: Part IV --- A structural Pandrosion
reduction for Smale's mean value conjecture}.
Preprint, April 2026.

\bibitem{paper5} I.~Besevic,
\emph{Universitas Pandrosion: Part V --- The quartic case of Smale's
conjecture in Pandrosion form}.
Preprint, April 2026.

\bibitem{smale1981} S.~Smale,
\emph{The fundamental theorem of algebra and complexity theory}.
Bull. Amer. Math. Soc.\ \textbf{4} (1981), 1--36.

\bibitem{beardon} A.~F.~Beardon, D.~Minda, T.~W.~Ng,
\emph{Smale's mean value conjecture and the hyperbolic metric}.
Math. Ann.\ \textbf{322} (2002), 623--632.

\bibitem{tischler} D.~Tischler,
\emph{Critical points and values of complex polynomials}.
J. Complexity \textbf{5} (1989), 438--456.

\bibitem{conte-vu} D.~Conte, V.~Vu,
\emph{Smale's mean value conjecture for finite Blaschke products}.
Comput. Methods Funct. Theory \textbf{8} (2008), 85--104.

\bibitem{dubinin} V.~N.~Dubinin, T.~Sugawa,
\emph{Geometric estimates for the Smale conjecture}.
J. Anal. Math.\ \textbf{114} (2011), 99--112.

\bibitem{lojasiewicz} S.~\L{}ojasiewicz,
\emph{Sur les ensembles semi-analytiques}.
Actes du Congrès International des Mathématiciens (Nice, 1970),
Tome~2, Gauthier-Villars, 1971, 237--241.

\bibitem{bochnak} J.~Bochnak, M.~Coste, M.-F.~Roy,
\emph{Real Algebraic Geometry}.
Springer, 1998.
\end{thebibliography}

\paragraph{Verification scripts.}
The companion scripts
\verb|/Users/ivanbesevic/Documents/poussiere/closure_d5_verify.py|
and
\verb|/Users/ivanbesevic/Documents/poussiere/d5_full_closure.py|
recompute every constant and verify every intermediate claim of this
paper from scratch, in sympy and numpy. They contain no dependency on
the existing scripts in \texttt{latex/scripts/}.

\end{document}
